\section{Lineare Abhängigkeit und Unabhängigkeit}
\label{sec:vector_algebra_linear_dependence}

Vektoren beschreiben Verschiebungen. Daher stellt sich die Frage, ob eine
Verschiebung durch eine gewichtete Kombination anderer Verschiebungen erzeugt
werden kann. In technischen Modellen zeigt dies beispielsweise, ob eine
zusätzliche Kraft- oder Stellgröße eine neue Wirkung ermöglicht oder durch
bereits vorhandene Größen ersetzt werden kann. Diese Frage führt auf die
lineare Abhängigkeit und Unabhängigkeit von Vektoren.

\subsection{Zusammensetzen und kompensieren}
\label{subsec:vector_algebra_linear_dependence_construction}

Wir beginnen mit der Frage, ob sich ein Vektor $\vec c$ aus zwei anderen
Vektoren zusammensetzen lässt:
\[
  \text{Gibt es }\alpha,\beta\in\R\text{ mit}
  \qquad
  \alpha\vec a+\beta\vec b=\vec c\,?
\]
Durch Verschieben von $\vec c$ auf die linke Seite erhalten wir die
gleichwertige Kompensationsfrage:
\[
  \alpha\vec a+\beta\vec b=\vec c
  \quad\Longleftrightarrow\quad
  \alpha\vec a+\beta\vec b-\vec c=\vec 0.
\]
Die von $\alpha\vec a$ und $\beta\vec b$ erzeugte Verschiebung kann also durch
$-\vec c$ vollständig kompensiert werden. Beide Sichtweisen beschreiben
denselben geschlossenen Vektorzug.

\begin{center}
  \CachedTikzFigure{content/resources/chapters/vector_algebra/tikz/linear_dependence_construction_and_compensation.tex}
\end{center}

Eine gewichtete Summe
$\lambda_1\vec v_1+\dots+\lambda_k\vec v_k$ heißt
\emph{Linearkombination} \EnglishTerm{linear combination} der Vektoren. Die
Ausgangsfrage lautet somit: Ist $\vec c$ eine Linearkombination von $\vec a$
und $\vec b$?

Allgemeiner suchen wir Koeffizienten $\alpha,\beta,\gamma$, für die
\[
  \alpha\vec a+\beta\vec b+\gamma\vec c=\vec 0
\]
gilt. In Koordinaten ist dies das homogene lineare Gleichungssystem
\[
  \begin{aligned}
    a_1\alpha+b_1\beta+c_1\gamma&=0,\\
    a_2\alpha+b_2\beta+c_2\gamma&=0,\\
    &\mathrel{\phantom{=}}\vdots\\
    a_n\alpha+b_n\beta+c_n\gamma&=0.
  \end{aligned}
  \qquad\Longleftrightarrow\qquad
  \left[
    \begin{array}{ccc|c}
      a_1&b_1&c_1&0\\
      a_2&b_2&c_2&0\\
      \vdots&\vdots&\vdots&\vdots\\
      a_n&b_n&c_n&0
    \end{array}
  \right].
\]
Die ursprüngliche Baufrage ist darin als Spezialfall $\gamma=-1$ enthalten.

\begin{DefinitionBox}{Lineare Abhängigkeit und Unabhängigkeit}{vector_algebra_linear_dependence}
  Die Vektoren $\vec v_1,\dots,\vec v_k\in\R^n$ heißen \emph{linear abhängig}
  \EnglishTerm{linearly dependent}, wenn das homogene Gleichungssystem
  \[
    \lambda_1\vec v_1+\dots+\lambda_k\vec v_k=\vec 0
  \]
  eine nichttriviale Lösung besitzt, also mindestens ein
  $\lambda_i\ne0$ ist. Besitzt es nur die triviale Lösung
  $\lambda_1=\dots=\lambda_k=0$, heißen die Vektoren \emph{linear unabhängig}
  \EnglishTerm{linearly independent}.

  Da ein homogenes LGS nicht ohne Lösung sein kann, ergibt sich damit die
  folgende Zuordnung:
  \begin{center}
    \CachedTikzFigure{content/resources/chapters/vector_algebra/tikz/linear_dependence_solution_types.tex}
  \end{center}
\end{DefinitionBox}

Damit sind die folgenden Fragen gleichwertig:
\begin{itemize}[leftmargin=1.7em,itemsep=0.45em]
  \item Lässt sich einer der Vektoren aus den anderen zusammensetzen?
  \item Können sich gewichtete Verschiebungen der Vektoren gegenseitig
    kompensieren, ohne dass alle Gewichte null sind?
  \item Besitzt das zugehörige homogene LGS eine nichttriviale Lösung?
\end{itemize}

\subsection{Entscheidung mithilfe der Determinante}
\label{subsec:vector_algebra_linear_dependence_determinant}

Nach Merke~\ref{key:linear_systems_homogeneous_solution_sets} besitzt ein
homogenes LGS entweder nur die triviale oder unendlich viele Lösungen. Für
genau $n$ Vektoren $\vec v_1,\dots,\vec v_n\in\R^n$ ist die zugehörige
Koeffizientenmatrix $M$ quadratisch. Ihre Spalten enthalten die Koordinaten der
Vektoren. Damit entscheidet die Determinante zwischen den beiden möglichen
Lösungsmengen.

Ist $\det(M)\ne0$, besitzt das LGS genau eine Lösung. Da die triviale Lösung
bei einem homogenen LGS immer vorhanden ist, muss sie die einzige Lösung sein.
Ist dagegen $\det(M)=0$, besitzt das LGS nicht genau eine Lösung. Da seine
Lösungsmenge auch nicht leer sein kann, bleiben nur unendlich viele Lösungen,
darunter auch nichttriviale.

\begin{TheoremBox}{Determinantenkriterium}{vector_algebra_linear_dependence_determinant_criterion}
  Für $n$ Vektoren $\vec v_1,\dots,\vec v_n\in\R^n$ und ihre
  Koeffizientenmatrix $M$ gilt
  \[
    \begin{aligned}
      \det(M)\ne0
      &\quad\Longleftrightarrow\quad
      \text{das homogene LGS besitzt nur die triviale Lösung}\\[-0.1em]
      &\quad\Longleftrightarrow\quad
      \vec v_1,\dots,\vec v_n\text{ sind }
      \Cloze{4.1cm}{\text{linear unabhängig}},\\[0.8em]
      \det(M)=0
      &\quad\Longleftrightarrow\quad
      \text{das homogene LGS besitzt nichttriviale Lösungen}\\[-0.1em]
      &\quad\Longleftrightarrow\quad
      \vec v_1,\dots,\vec v_n\text{ sind }
      \Cloze{3.7cm}{\text{linear abhängig}}.
    \end{aligned}
  \]
\end{TheoremBox}

\begin{ExampleBox}{Unabhängigkeit mit der Determinante prüfen}{vector_algebra_linear_independence_determinant}
  Für
  \[
    \vec a=\begin{pmatrix}2\\1\\1\end{pmatrix},\qquad
    \vec b=\begin{pmatrix}2\\3\\4\end{pmatrix},\qquad
    \vec c=\begin{pmatrix}7\\9\\11\end{pmatrix}
  \]
  prüfen wir die Koeffizientenmatrix mit den Vektoren als Spalten:
  \[
    \det(M)
    =\begin{ClozeContinuation}[25]{5}
      \det\!\begin{pmatrix}
        2&2&7\\
        1&3&9\\
        1&4&11
      \end{pmatrix}
      =2(33-36)-2(11-9)+7(4-3)
      =-3
    \end{ClozeContinuation}.
  \]
  Da $\det(M)\ne0$ ist, sind $\vec a$, $\vec b$ und $\vec c$
  \Cloze{4.1cm}{linear unabhängig}.
\end{ExampleBox}

\begin{RemarkBox}[title=Wann ist die Determinante anwendbar?]
  Das Determinantenkriterium gilt unmittelbar für $n$ Vektoren im $\R^n$.
  Bei einer anderen Anzahl von Vektoren untersuchen wir das homogene LGS mit
  dem Gauß-Verfahren. Mehr als $n$ Vektoren im $\R^n$ sind stets linear
  abhängig.
\end{RemarkBox}

\subsection{Geometrische Bedeutung}
\label{subsec:vector_algebra_linear_dependence_geometry}

In zwei und drei Dimensionen lässt sich die Determinante geometrisch deuten.
Für drei Vektoren im $\R^3$ ist sie gerade das Spatprodukt:
\[
  \det(\vec a\ \vec b\ \vec c)
  =(\vec a\times\vec b)\cdot\vec c.
\]

\begin{KeyPointBox}{Fläche und Volumen zeigen Unabhängigkeit}{vector_algebra_linear_dependence_geometry}
  Unabhängige Richtungen spannen einen nicht verschwindenden geometrischen
  Inhalt auf:
  \[
    \begin{aligned}
      \vec a,\vec b\in\R^2\text{ linear unabhängig}
      &\quad\Longleftrightarrow\quad
      |\det(\vec a\ \vec b)|>0
      &&\text{(Flächeninhalt)},\\[0.7em]
      \vec a,\vec b,\vec c\in\R^3\text{ linear unabhängig}
      &\quad\Longleftrightarrow\quad
      |(\vec a\times\vec b)\cdot\vec c|>0
      &&\text{(Volumen)}.
    \end{aligned}
  \]
  Kollabiert die aufgespannte Fläche oder das aufgespannte Volumen zu null,
  sind die beteiligten Vektoren linear abhängig.
\end{KeyPointBox}
